\begin{textsample}{POS Dim 2 – rr – Score 25.00 – rr/rr\_ef\_50.txt}  \label{ex:f2_pos_149}
4.12 MODALIDADE DE EDUCAÇÃO DE JOVENS E ADULTOS – EJA Pensar em um \textbf{currículo} direcionado para a Educação de Jovens e Adultos – EJA, em Roraima, requer um conhecimento do perfil deste aluno, que em sua maioria é um \textbf{trabalhador} ou dona de casa, que está em busca de uma melhoria socioeconômica, cultural e que possibilite o exercício pleno de sua cidadania.
A modalidade EJA \textbf{ofertada} na Rede Pública Estadual de Roraima de forma presencial com avaliação no processo, considerando a idade mínima de 15 anos completos para a matrícula no Ensino Fundamental, e de 18 ( dezoito ) anos completos, no Ensino Médio, para os alunos que não tiveram acesso ou continuidade de estudos em idade própria.
A \textbf{carga} \textbf{horária} no Ensino Fundamental está distribuída da seguinte forma : 1º segmento ( 2º ao 5º ano ) 1.600 horas e 2º segmento ( 6º ao 9º ano ) 2.000 horas.
A Declaração de Hamburgo ( 1997 ) ressalta que é objetivo da Educação de Jovens e Adultos o desenvolvimento da autonomia e do senso de responsabilidade das pessoas e das comunidades, sendo, portanto necessário um modelo de educação onde o diálogo seja prioridade entre as partes.
É essencial que os enfoques da educação de adultos estejam baseados no patrimônio, na cultura, nos valores e nas experiências anteriores das pessoas, e que as distintas maneiras de por em prática estes enfoques facilitem e estimulem a ativa participação e expressão do educando. ( Declaração de Hamburgo sobre Educação de Adultos.
Item 5, 1997 ) A Lei de \textbf{Diretrizes} e Bases da Educação – LDB nº 9.394/96, inclui a EJA como modalidade da Educação Básica e dedica especificamente dois \textbf{artigos}, 37 e 38 que reafirmam a obrigatoriedade, a gratuidade da \textbf{oferta} da educação para todos e define como público alvo, àqueles que não tiveram acesso ou continuidade de estudos no Ensino Fundamental e Médio na idade própria.
No Estado de Roraima, a EJA teve início em 1948, no prédio da Escola Lobo D'Almada, com o nome de Escola Supletiva Noturna. \textbf{Atendia} alunos de 1ª a 4ª \textbf{série} do \textbf{Curso} Primário, como era denominada.
A partir de 1951, passou a chamar -se “ \textbf{Curso} Lourenço Filho ” e funcionou no mesmo prédio, até 1972.
Com o advento da Lei 5.692/71 ( Lei de \textbf{Diretrizes} e Bases ), o sistema supletivo de ensino foi implantado em todo país, em forma de \textbf{cursos} e \textbf{exames}, \textbf{atendendo} a um contingente cada vez maior de adultos, a cada ano.
Em Roraima, os \textbf{Exames} de Suplência Geral começaram a ser realizados em 1972.
Também foram realizados \textbf{Exames} de Suplência Profissionalizante, voltados para uma \textbf{clientela} inserida na força de trabalho especializado.
O primeiro \textbf{curso} voltado para a \textbf{clientela} analfabeta foi o Movimento Brasileiro de Alfabetização - MOBRAL.
Os programas posteriores procuravam dar continuidade aos estudos dos recém-alfabetizados, com precário domínio da leitura e escrita.
Destacamos ainda, o Programa de Educação Integrada - PEI, realizado em parceria entre o Governo Federal e a Secretaria de Estadual Educação, funcionava com aulas presenciais, no turno noturno, no nível de 1ª a 4ª \textbf{série}, com duração de um ano letivo.
A nomenclatura foi alterada posteriormente passando a denominar -se Programa de Educação Básica – PEB, com os mesmos objetivos.
Para oportunizar a continuidade dos estudos ao público alvo, foi criado, em 1975, o primeiro Centro de Estudos Supletivos - CES, conhecido como CES- São Francisco, no \textbf{município} de Boa Vista.
O Centro \textbf{atendeu}, inicialmente, a \textbf{clientela} maior de 18 anos, no nível de 5ª a 8ª \textbf{série} do Ensino Fundamental.
Nos \textbf{cursos} profissionalizantes de Ensino Médio, podemos destacar o PROJETO LOGOS II que era \textbf{destinado} a habilitar professores leigos, teve sua atuação nos \textbf{municípios} de Boa Vista, Caracaraí, Normandia, Alto Alegre e São Luis, entre os anos de 1976 a 1990, obedecendo o sistema de ensino por meio de módulos.
Outra ação desenvolvida foi o Projeto Minerva, \textbf{Curso} de Suplência Geral no nível de 5ª a 8ª \textbf{série}, veiculado pelo rádio para todo o Estado.
O Projeto Minerva foi substituído pelo Supletivo de 1º Grau - SPG, que funcionava nas escolas estaduais, nos mesmos moldes do anterior, sendo que as aulas eram gravadas e ouvidas pelos alunos através de gravador portátil.
Com a instalação da Rede Amazônica de Televisão, em Boa Vista, foi veiculado o primeiro \textbf{curso} supletivo televisivo – o Projeto João da Silva, que \textbf{atendeu} alunos de 1ª a 4ª \textbf{série}.
Em 1990 foi implantado na Penitenciária Agrícola de Monte Cristo, o programa SPG.
A partir de 2001, foi criado um \textbf{curso} preparatório para os \textbf{exames} de Suplência Geral, nos níveis de ensino fundamental e médio, desenvolvido através da Educação a Distância, utilizando material do Telecurso 2000, tendo sido interrompido, no ano de 2007.
Posteriormente outras experiências foram vivenciadas, como por exemplo, o Programa Estadual de Teleducação para o Ensino Médio - PETEM e o Programa Estadual de Teleducação para o Ensino Fundamental – PETEF, aliando, tele aulas, material de apoio impresso e monitores treinados para acompanhar a aprendizagem dos alunos.
Em 1995 foi criado e implantado nas Escolas Estaduais o \textbf{Curso} Seriado Supletivo, que oportunizava aos que não tiveram acesso ou continuidade de estudos no tempo certo, em regime de aula presencial, o aluno fazia duas \textbf{séries} por ano, concluindo a etapa do Ensino Fundamental.
Em 1998, o \textbf{Curso} Seriado Supletivo foi substituído pelo Projeto Piracema, que também trabalhava com aulas em regimes presenciais, porém com o diferencial de oportunizar a conclusão de \textbf{séries}, a cada \textbf{semestre} concluído com êxito.
Em 2002 foi aprovada pelo CEE/RR a Proposta de Reestruturação dos \textbf{Cursos} da EJA, através do \textbf{Parecer} CEE/RR nº 59/02, sendo o atendimento realizado de forma presencial, com avaliação no processo, em \textbf{cursos} semestrais, nas escolas localizadas em todo o Estado.
Em 2004, foi oficializado o Fórum Permanente de Debates da Educação de Jovens e Adultos do Estado de Roraima, com a participação da sociedade civil, para discutir as \textbf{políticas} públicas da EJA.
A \textbf{Legislação} em relação EJA, apresentam “ avanços ” educacionais significativos no decorrer do processo, principalmente em relação a LDB, conforme descrito na tabela abaixo : Tabela 1 : Evolução da EJA nas Leis de \textbf{Diretrizes} e Bases da Educação Nacional

% matched lemmas: artigo, atender, carga, clientela, currículo, curso, destinar, diretriz, exame, horário, legislação, município, oferta, ofertar, parecer, política, semestre, série, trabalhador
\end{textsample}
